% GENERATED FILE. Edit canonical lessons, then run npm run generate:books.
% canonical-source-hash: fnv1a64:53b9e1b749029754
% canonical-lessons: JA-C393-mejirushi, JA-C393-kanban, JA-C393-hyoushiki, JA-C393-yajirushi, JA-C393-furimuku

\chapter{Signs and Landmarks}
\label{ch:ja-signs-and-landmarks}
\hlchaptermodality{393}

\begin{chapteropening}
\textbf{By the end of this chapter:} \emph{I can follow signs and landmarks.}

Complete the last lesson of chapter 393: I can follow signs and landmarks.
\end{chapteropening}

\section[\texorpdfstring{\ja{めじるし} (mejirushi)}{めじるし (mejirushi)}]{\ja{めじるし} (mejirushi) — a landmark; a mark to go by}

\label{lesson:JA-C393-mejirushi}



\begin{quote}
Before the new one: say the Japanese for to turn back, then the Japanese for to stop by.
\end{quote}

\subsection*{You'll want to know: \ja{めじるし}}
\textbf{\ja{めじるし}} — \emph{mejirushi} — \textquotedblleft{}a landmark; a mark to go by\textquotedblright{}.

\begin{quote}
\textbf{\ja{えきの} \ja{まえの} \ja{たかい} \ja{たてものが} \ja{めじるしです。}} — \emph{eki no mae no takai tatemono ga mejirushi desu} — \textquotedblleft{}The tall building in front of the station is the landmark.\textquotedblright{}
\end{quote}

\begin{cousinweb}
\textbf{\ja{めじるし}} (\emph{mejirushi}) is a native Japanese word (wago), inherited rather than borrowed from Chinese, built from \emph{me}, \textquotedblleft{}eye\textquotedblright{} and \emph{jirushi}, \textquotedblleft{}mark\textquotedblright{}; its kanji were attached later as kun'yomi (JMdict).
\end{cousinweb}

\begin{grammarlens}[title={What you've built}]
One more word for getting around.
\end{grammarlens}

\subsection*{Your turn}
\begin{itemize}
\raggedright
  \item \emph{Say it:} \emph{mejirushi}
  \item \emph{Say it:} \emph{mejirushi}, once more
  \item \emph{Say it:} \emph{eki no mae no takai tatemono ga mejirushi desu} — The tall building in front of the station is the landmark.
  \item \emph{Recall:} say the Japanese for to turn back, then the Japanese for to stop by, then say \emph{mejirushi} again
\end{itemize}

\begin{tcolorbox}[breakable,colback=teal!4,colframe=teal!35!black,title={Before you move on}]
What does \ja{めじるし} mean? (\textquotedblleft{}A landmark; a mark to go by\textquotedblright{}.) Say it once more.
\end{tcolorbox}
\section[\texorpdfstring{\ja{かんばん} (kanban)}{かんばん (kanban)}]{\ja{かんばん} (kanban) — a signboard; a shop sign}

\label{lesson:JA-C393-kanban}



\begin{quote}
Before the new one: say the Japanese for to stop by, then the Japanese for a landmark.
\end{quote}

\subsection*{You'll want to know: \ja{かんばん}}
\textbf{\ja{かんばん}} — \emph{kanban} — \textquotedblleft{}a signboard; a shop sign\textquotedblright{}.

\begin{quote}
\textbf{\ja{みせの} \ja{かんばんは} \ja{あかいです。}} — \emph{mise no kanban wa akai desu} — \textquotedblleft{}The shop's sign is red.\textquotedblright{}
\end{quote}

\begin{cousinweb}
\textbf{\ja{かんばん}} (\emph{kanban}) is Sino-Japanese (kango): two Chinese characters read with their on'yomi, readings borrowed from Chinese. They are \emph{kan}, \textquotedblleft{}watch over\textquotedblright{} (Middle Chinese \emph{khan}) and \emph{ban}, \textquotedblleft{}board\textquotedblright{} (\emph{paenX}) (KANJIDIC2; Middle Chinese in Baxter's transcription).
\end{cousinweb}

\begin{grammarlens}[title={What you've built}]
One more word for getting around.
\end{grammarlens}

\subsection*{Your turn}
\begin{itemize}
\raggedright
  \item \emph{Say it:} \emph{kanban}
  \item \emph{Say it:} \emph{kanban}, once more
  \item \emph{Say it:} \emph{mise no kanban wa akai desu} — The shop's sign is red.
  \item \emph{Recall:} say the Japanese for to stop by, then the Japanese for a landmark, then say \emph{kanban} again
\end{itemize}

\begin{tcolorbox}[breakable,colback=teal!4,colframe=teal!35!black,title={Before you move on}]
What does \ja{かんばん} mean? (\textquotedblleft{}A signboard; a shop sign\textquotedblright{}.) Say it once more.
\end{tcolorbox}
\section[\texorpdfstring{\ja{ひょうしき} (hyōshiki)}{ひょうしき (hyōshiki)}]{\ja{ひょうしき} (hyōshiki) — a road sign; a marker}

\label{lesson:JA-C393-hyoushiki}



\begin{quote}
Before the new one: say the Japanese for a landmark, then the Japanese for a signboard.
\end{quote}

\subsection*{You'll want to know: \ja{ひょうしき}}
\textbf{\ja{ひょうしき}} — \emph{hyōshiki} — \textquotedblleft{}a road sign; a marker\textquotedblright{}.

\begin{quote}
\textbf{\ja{ひょうしきに} \ja{ちゅういして} \ja{ください。}} — \emph{hyōshiki ni chūi shite kudasai} — \textquotedblleft{}Please pay attention to the road signs.\textquotedblright{}
\end{quote}

\begin{cousinweb}
\textbf{\ja{ひょうしき}} (\emph{hyōshiki}) is Sino-Japanese (kango): two Chinese characters read with their on'yomi, readings borrowed from Chinese. They are \emph{hyō}, \textquotedblleft{}signpost\textquotedblright{} (Middle Chinese \emph{pjiew}) and \emph{shiki}, \textquotedblleft{}know\textquotedblright{} (\emph{syik}) (KANJIDIC2; Middle Chinese in Baxter's transcription).
\end{cousinweb}

\begin{grammarlens}[title={What you've built}]
One more word for getting around.
\end{grammarlens}

\subsection*{Your turn}
\begin{itemize}
\raggedright
  \item \emph{Say it:} \emph{hyōshiki}
  \item \emph{Say it:} \emph{hyōshiki}, once more
  \item \emph{Say it:} \emph{hyōshiki ni chūi shite kudasai} — Please pay attention to the road signs.
  \item \emph{Recall:} say the Japanese for a landmark, then the Japanese for a signboard, then say \emph{hyōshiki} again
\end{itemize}

\begin{tcolorbox}[breakable,colback=teal!4,colframe=teal!35!black,title={Before you move on}]
What does \ja{ひょうしき} mean? (\textquotedblleft{}A road sign; a marker\textquotedblright{}.) Say it once more.
\end{tcolorbox}
\section[\texorpdfstring{\ja{やじるし} (yajirushi)}{やじるし (yajirushi)}]{\ja{やじるし} (yajirushi) — an arrow (on a sign)}

\label{lesson:JA-C393-yajirushi}



\begin{quote}
Before the new one: say the Japanese for a signboard, then the Japanese for a road sign.
\end{quote}

\subsection*{You'll want to know: \ja{やじるし}}
\textbf{\ja{やじるし}} — \emph{yajirushi} — \textquotedblleft{}an arrow (on a sign)\textquotedblright{}.

\begin{quote}
\textbf{\ja{やじるしの} \ja{ほうに} \ja{あるいて} \ja{ください。}} — \emph{yajirushi no hō ni aruite kudasai} — \textquotedblleft{}Please walk the way the arrow points.\textquotedblright{}
\end{quote}

\begin{cousinweb}
\textbf{\ja{やじるし}} (\emph{yajirushi}) is a native Japanese word (wago), inherited rather than borrowed from Chinese, built from \emph{ya}, \textquotedblleft{}arrow\textquotedblright{} and \emph{jirushi}, \textquotedblleft{}mark\textquotedblright{}; its kanji were attached later as kun'yomi (JMdict).
\end{cousinweb}

\begin{grammarlens}[title={What you've built}]
One more word for getting around.
\end{grammarlens}

\subsection*{Your turn}
\begin{itemize}
\raggedright
  \item \emph{Say it:} \emph{yajirushi}
  \item \emph{Say it:} \emph{yajirushi}, once more
  \item \emph{Say it:} \emph{yajirushi no hō ni aruite kudasai} — Please walk the way the arrow points.
  \item \emph{Recall:} say the Japanese for a signboard, then the Japanese for a road sign, then say \emph{yajirushi} again
\end{itemize}

\begin{tcolorbox}[breakable,colback=teal!4,colframe=teal!35!black,title={Before you move on}]
What does \ja{やじるし} mean? (\textquotedblleft{}An arrow (on a sign)\textquotedblright{}.) Say it once more.
\end{tcolorbox}
\section[\texorpdfstring{\ja{ふりむく} (furimuku)}{ふりむく (furimuku)}]{\ja{ふりむく} (furimuku) — to turn round; to look back}

\label{lesson:JA-C393-furimuku}



\begin{quote}
Before the new one: say the Japanese for a road sign, then the Japanese for an arrow (on a sign).
\end{quote}

\subsection*{You'll want to know: \ja{ふりむく}}
\textbf{\ja{ふりむく}} — \emph{furimuku} — \textquotedblleft{}to turn round; to look back\textquotedblright{}.

\begin{quote}
\textbf{\ja{うしろから} \ja{こえが} \ja{したので、ふりむきました。}} — \emph{ushiro kara koe ga shita node, furimukimashita} — \textquotedblleft{}I heard a voice behind me, so I turned round.\textquotedblright{}
\end{quote}

Polite \textbf{\ja{ふりむきます}}, \emph{furimukimasu}; a godan verb. Past (polite): \textbf{\ja{ともだちが} \ja{ふりむきました。}}, \emph{tomodachi ga furimukimashita}.

\begin{cousinweb}
\textbf{\ja{ふりむく}} (\emph{furimuku}) is a native Japanese word (wago), inherited rather than borrowed from Chinese; its kanji was attached later as a kun'yomi (JMdict).
\end{cousinweb}

\begin{grammarlens}[title={What you've built}]
One more word for getting around.
\end{grammarlens}

\subsection*{Your turn}
\begin{itemize}
\raggedright
  \item \emph{Say it:} \emph{furimuku}
  \item \emph{Say it:} \emph{furimuku}, once more
  \item \emph{Say it:} \emph{ushiro kara koe ga shita node, furimukimashita} — I heard a voice behind me, so I turned round.
  \item \emph{Recall:} say the Japanese for a road sign, then the Japanese for an arrow (on a sign), then say \emph{furimuku} again
\end{itemize}

\begin{tcolorbox}[breakable,colback=teal!4,colframe=teal!35!black,title={Before you move on}]
What does \ja{ふりむく} mean? (\textquotedblleft{}To turn round; to look back\textquotedblright{}.) Say it once more.
\end{tcolorbox}
